\documentclass[11pt]{article}
\usepackage[a4paper,margin=27mm]{geometry}
\usepackage{amsmath,amssymb,amsthm,hyperref}
\newtheorem{theorem}{Theorem}
\newtheorem{lemma}{Lemma}
\newcommand{\E}{\mathbb E}
\newcommand{\Pol}{\operatorname{Pol}}
\title{An elementary endpoint distribution-function bootstrap}
\author{Dependency simplification: independent mathematical reviews completed}
\date{30 September 2026}
\begin{document}\maketitle
\noindent\textbf{Status.} Independently mathematically reviewed, with no
gap identified; not formally verified.
This note replaces the classical localized-Laguerre input previously
used for the exponent-$3/4$ distribution-function estimate. It uses
only the reviewed Gaussian comparison and subset estimates. It does
not assume the quadratic per-die theorem.

Use the notation of \path{../power_per_die_lower_bound.tex}:
$u_j=1-p_j$, $s=\sum_j u_j$, $\ell=\lfloor m/4\rfloor$,
$N=m-\ell$, $\alpha=\ell/m$, $x=\alpha s$, and $\theta=1/16$.
For a uniform $\ell$-subset $B$, write
$A_B=\prod_{j\in B}p_j$, $\bar A=\E_B A_B$,
and $v_B=N^{-1}\sum_{j\notin B}\ell u_j$.
Let $\mu$ have row masses $(\ell+1)w\bar A$, let $\nu$ have row
masses $mw$, and let $\beta$ have density
$b_\ell(z)=(\ell+1)\ell^{-1}(1-z/\ell)^\ell$ on $[0,\ell]$.
All constants are absolute.

\begin{theorem}\label{main}
For every monotone tensor with the exact uniform mixed moments,
\[
 \nu([0,t])\le C(t+m^{-3/4})\qquad(t\ge0).
\]
\end{theorem}

The reviewed comparison supplies a starting estimate
\begin{equation}\label{base}
 \nu([0,t])\le C(t+h),\qquad h=m^{-1/2}.
\end{equation}
Here is a direct derivation, to specify the input precisely. Put
\[
 T_d(y)=\sum_{j<d}(1-j/d)L_j(y)=d^{-1}L_{d-1}^{[2]}(y),
 \qquad U_d(z)=d^{-2}T_d(\theta z)^2.
\]
Its Laguerre coefficient product is bounded by an absolute constant.
The reviewed bilinear comparison gives
$\mu(U_d)\le\beta(U_d)+Cd/m\le C/d+Cd/m$ for $d\le\sqrt m$.
For $0\le x\le c_0/d$, $U_d(x)\ge c>0$ by the elementary
near-zero Laguerre estimate. Choose $c_0$ so small that $x\le c_0$
implies $s\le1/2$; then $\bar A\ge1-s\ge1/2$ and $\nu$ is bounded
by a constant times $\mu$ there. Taking
$d\asymp\min(\sqrt m,1/t)$ proves \eqref{base} near zero. For larger
$t$, the finite row-mass estimate $\nu([0,t])\le C(t+1)$ suffices.

\begin{lemma}\label{comparison}
Uniformly for $m^{1/4}\le d\le m^{3/4}$,
\[
 |\mu(U_d)-\beta(U_d)|
 \le C\left(\frac1m+\frac{hd}m+\frac{\log d}{dm}\right)+O(m^{-3}).
\]
\end{lemma}
\begin{proof}
The Gaussian off-diagonal inequality states
$y^{a/2}|L_r^{[a]}(y)|\le e^{y/2}\sqrt{(r+a)!/r!}$.
Since
\[
 D_z^j T_d(\theta z)=(-\theta)^jL_{d-j-1}^{[j+2]}(\theta z)/d,
\]
the total matrix index is always $d+1$. It follows, for every
derivative order $j$ simultaneously, that
\[
 |D_z^jT_d(\theta z)|\le C e^{\theta z/2}z^{-1}
                                  (C\sqrt{d/z})^j.
\]
Thus, for all $z>0$ and all $k$,
\begin{equation}\label{Uder}
 |U_d^{(k)}(z)|/k!
       \le\frac{C e^{\theta z}}{d^2z^2}
                                \frac{(C\sqrt{d/z})^k}{k!}.
\end{equation}
The constants are uniform in $k$. The elementary whole-series bound
\begin{equation}\label{series}
 \sum_{k\ge2}(k+1)a^k k^{k/2}/k!\le Ca^2e^{Ca^2}
\end{equation}
follows by splitting even and odd orders and comparing to exponential
series; for example $(2j)^j/(2j)!\le2^j/j!$.

We use only the older balance $\max u_j\le C(s+1)/m$.
The centered variance of $z_j=\ell u_j$, $j\notin B$, is at most
$CN(s+1)v_B$. The centered elementary-coefficient bound therefore
cancels the factors $v_B^{-k/2}$ in \eqref{Uder}.
Under the success-weighted subset law, the reviewed bias and second
moment estimates are
\begin{equation}\label{moments}
 0\le\E\Delta\le Cx(x+1)/m,\qquad
 \E\Delta^2\le C\{x(x+1)/m+x^2(x+1)^2/m^2\},
 \quad\Delta=v_B-x,
\end{equation}
when $s\le m/256$. Also $\bar A\le e^{-x}$.

For $x\le1/d$, use the global Gaussian derivative bounds with
coefficient product $O(1)$. The complete polarization series and
the singular Taylor averaging estimate each have pointwise error
$O(d/m)$, because $s+1=O(1)$. They vanish if $x=0$.
Their integrated contribution, by \eqref{base}, is
\begin{equation}\label{small}
 O((d/m)(1/d+h))=O(1/m+hd/m).
\end{equation}

For $1/d<x\le d$, first restrict to subsets with $v_B\ge x/2$.
We always have $v_B<2x$. The variance estimate, \eqref{Uder}, and
\eqref{series} bound the complete polarization error by
\[
 \frac{C e^{\theta v_B}}{d^2x^2}
       \frac{d(s+1)}m\exp[Cd(s+1)/m].
\]
After averaging with $A_B$, $\bar A\le e^{-x}$ leaves a fixed
exponential margin, since $d/m\le m^{-1/4}\to0$.
The resulting envelope is
\begin{equation}\label{envelope}
 \frac C{dm}e^{-cx}(x^{-1}+x^{-2}).
\end{equation}
For averaging, retain the full linear bias and split only the Taylor
remainder into good and bad subsets. On good subsets, use
$|U_d'|\le C d^{-3/2}x^{-5/2}e^{2\theta x}$,
$|U_d''|\le C d^{-1}x^{-3}e^{2\theta x}$, and \eqref{moments}.
The same envelope results: the ratio of the linear term to the
leading quadratic term is at most $C\sqrt{x/d}\le C$, and the
squared-bias term is absorbed in the exponential margin.

Here the bad-subset contribution is negligible. The uniform law
satisfies
\[
 \Pr_B(v_B<x/2)\le e^{-cm\min(x,1)}.
\]
Indeed the entire population $z_j$ has mean $x$ and maximum
$R=C(s+1)$; Maclaurin's inequality gives
$\E e^{-Nv_B/R}\le\exp[-Nx(1-e^{-1})/R]$, and Markov's inequality
gives the claim. The full Gaussian series and the coarse Taylor
remainder are polynomial in $m$ times $e^{Cx}$ on $x\le d$.
For $1/d<x\le1$, the bad exponent is at most $-cm/d$; for
$1\le x\le d$, it is at most $-cm+Cd$. Both are superpolynomially
small uniformly for $m^{1/4}\le d\le m^{3/4}$.
Weighted normalization causes no loss, since
$\bar A\Pr_{\rm weighted}(\mathrm{bad})
=\E_{\rm uniform}A_B\boldsymbol1_{\rm bad}
\le\Pr_{\rm uniform}(\mathrm{bad})$.

For $x>d$, the global Gaussian series, $\bar A\le e^{-x}$, and
$d/m=o(1)$ give a polynomial factor times $e^{-cx}$.
The nonsingular first derivative and uniform-subset Cauchy--Schwarz
give the same estimate for averaging. These tails are $O(m^{-3})$
because $d\ge m^{1/4}$.

Finally \eqref{base} and Stieltjes integration by parts imply
\[
 \int_{1/d}^\infty e^{-cx}x^{-1}\,d\nu\le C(\log d+hd),\qquad
 \int_{1/d}^\infty e^{-cx}x^{-2}\,d\nu\le C(d+hd^2).
\]
Possible atoms at the cutoff are included by the same CDF bound.
Integrate \eqref{envelope}, add \eqref{small} and the negligible tails,
and use the exact identity
$\beta(P)=(\ell+1)\E_w\E_B[A_B\Pol_N(P)]$.
This proves the lemma.
\end{proof}

\begin{proof}[Proof of Theorem~\ref{main}]
For every $m^{1/4}\le d\le m^{3/4}$, Lemma~\ref{comparison} gives
\[
 \mu(U_d)\le C/d+C\{1/m+m^{-1/2}d/m+\log d/(dm)\}
                   +O(m^{-3})\le C/d.
\]
The last inequality uses $d^2\le m^{3/2}$.
Near-zero positivity and comparison of $\mu$ with $\nu$, as in the
derivation of \eqref{base}, give $\nu([0,c_0/d])\le C/d$.
For $t\lesssim m^{-1/4}$ choose
$d\asymp\min(m^{3/4},1/t)$; for larger $t$ use \eqref{base}.
This proves $\nu([0,t])\le C(t+m^{-3/4})$ uniformly.
\end{proof}
\end{document}
